\section{Entidades, atributos e identificadores}

Uma entidade representa um assunto sobre o qual o sistema precisa registrar informações.
O identificador distingue uma ocorrência das demais. A primeira linha é apenas um exemplo.

\begin{table}[!ht]
\centering
\small
\setlength{\tabcolsep}{4pt}
\renewcommand{\arraystretch}{1.75}
\begin{tabularx}{\linewidth}{|>{\raggedright\arraybackslash}X
                              |>{\raggedright\arraybackslash}X
                              |>{\raggedright\arraybackslash}X
                              |>{\centering\arraybackslash}X|}
\hline
\textbf{Entidade} &
\textbf{Atributos necessários} &
\textbf{Identificador} &
\textbf{RF relacionado(s)} \\
\hline
\textbf{Solicitação} &
Número, data de abertura, situação &
Número da solicitação &
RF01 (ilustrativo) \\
\hline
 & & & \\ \hline
 & & & \\ \hline
 & & & \\ \hline
 & & & \\ \hline
 & & & \\ \hline
 & & & \\ \hline
 & & & \\ \hline
 & & & \\ \hline
 & & & \\ \hline
\end{tabularx}
\end{table}

\begin{warnbox}[Observação]
Se faltar linha, continue registrando as entidades na parte de trás da folha ou em uma folha adicional.
\end{warnbox}
